% =============================================================================
% related_work.tex -- "Benchmark Contamination" + detection + memorization/privacy.
% Owner: Subagent A (Literature Review Writer). All \cite keys -> ../references.bib.
%
% CONSISTENCY CONTRACT: every detection method described here is implemented and
% evaluated (detectors/), and every implemented method is described here. No method
% is named that is not tested, and vice versa (see docs/method_selection_memo.md).
%
% DIVERGENCES FROM THE PASTED DRAFT (flagged for review):
% * Typology standardized to verbatim / paraphrased / semantic (the project spine);
% the draft's "input-label" case is folded in as a severity axis.
% * Replaced Skywork/Wei-2023 definitional cites with primary sources + the survey.
% * Detection subsection now matches the IMPLEMENTED+RUN shortlist exactly: LOSS,
% Min-K%, Min-K%++, zlib, n-gram overlap, and the Oren permutation test. Guided
% prompting and neighbourhood/reference MIA are explicitly framed as related
% approaches we DO NOT evaluate (they are not in the implemented set).
% =============================================================================

\section{Related Work: Contamination, Memorization, and Privacy Leakage}
\label{sec:relatedwork}

\subsection{Defining benchmark contamination}
We adopt the standard definition: \emph{benchmark contamination} is the presence of
evaluation data, inputs, labels, or accompanying metadata, within a model's
pre-training corpus~\cite{golchin2024timetravel}. Contamination matters for two
reasons that this paper treats as inseparable. First, it invalidates evaluation: a
contaminated score conflates capability with retrieval, so the metric no longer
estimates generalization. Second, and central to our thesis, contamination is a
\emph{symptom of, and a measurable proxy for, unintended memorization}, and
memorization of evaluation data sits on the same mechanism that leaks sensitive
content from the corpus. We make this contamination~$\rightarrow$~memorization~$
\rightarrow$~leakage chain the object of empirical study.

\subsection{A typology of contamination}
Following the project's framing and the contamination-detection
survey~\cite{ravaut2024survey}, we distinguish three forms by the transformation
between the corpus copy and the benchmark item:
\begin{itemize}
 \item \textbf{Verbatim contamination.} The exact token sequence of a test item
 appears in training data. This is what classical $n$-gram decontamination targets
 (e.g., the 13-gram overlap test introduced for GPT-3~\cite{brown2020gpt3}) and what
 verbatim-extraction memorization measures~\cite{carlini2023quantifying}.
 \item \textbf{Paraphrased contamination.} The semantic content is present but
 reworded, so surface-level $n$-gram matching misses it. A perfect verbatim filter
 provides only a false sense of safety, since style-transfer rephrasings evade it
 while preserving the leaked information~\cite{ippolito2023verbatim}.
 \item \textbf{Semantic contamination.} The underlying knowledge or answer is encoded
 without lexical overlap (e.g., the same question-answer mapping in a different
 format). Detecting it requires model-behavioral or distributional signals rather than
 string matching.
\end{itemize}
A second, orthogonal severity axis is \emph{what} is contaminated: input-only leakage
inflates familiarity, whereas joint input--label leakage enables direct answer
retrieval and is the most damaging to evaluation validity. Empirically, overlap between
open-model training data and benchmarks such as GSM8K has been reported for models
trained on largely undisclosed corpora~\cite{touvron2023llama}, motivating
ground-truth-controlled study on models whose corpus is fully public.

\subsection{Why memorization is a security and privacy problem}
Memorization is not a benign curiosity. Over-parameterized models trained on
web-scale scrapes retain and can regurgitate verbatim sequences, including personally
identifiable information (PII) such as names, emails, and phone
numbers~\cite{carlini2021extracting}. This has been formalized along several axes that
we reuse as outcome variables:
\begin{itemize}
 \item \textbf{$k$-eidetic / extractable memorization.} A string is extractable if a
 prefix makes the model regenerate it, and is $k$-eidetic if it occurs in at most $k$
 training documents~\cite{carlini2021extracting}; the prefix-continuation form under
 greedy decoding makes this directly measurable~\cite{carlini2023quantifying}.
 \item \textbf{Exposure and example-level memorization.} Injecting a canary secret and
 measuring its \emph{exposure}, its rank against random alternatives, quantifies
 unintended memorization and its growth with occurrence count~\cite{carlini2019secret};
 this requires control over the training process (canary insertion), which our
 pretrained-checkpoint setting does not afford, so we use it for definitions rather than
 as a measurement. Relatedly, memorization is concentrated on specific
 examples~\cite{zhang2023counterfactual} rather than spread uniformly, which is what
 makes per-item contamination scores meaningful predictors of per-item leakage.
 \item \textbf{Extraction at scale.} Production models can be driven, via a divergence
 attack, to emit memorized training data well above their nominal aligned rate,
 recovering thousands of verbatim examples cheaply~\cite{nasr2025scalable}.
 \item \textbf{PII leakage games.} Leakage of personally identifiable information
 decomposes into extraction, reconstruction, and inference; data scrubbing and
 differential privacy reduce but do not eliminate it~\cite{lukas2023pii}, models leak
 PII through memorization more than through associative
 inference~\cite{huang2022leaking}, and black-box probing tools can elicit a data
 subject's PII directly from a deployed model~\cite{kim2023propile}.
\end{itemize}
The security framing follows directly: if contamination is a measurable proxy for
memorization, and memorization is the vector for PII and proprietary-data exposure,
then contamination is not only a metrics problem but a \emph{privacy vulnerability}.

\subsection{The membership-inference lineage}
\label{sec:mia-lineage}
Deciding whether a specific record was in a model's training set is the canonical
privacy attack, and contamination detection is an instance of it. The lineage we build
on runs as follows. \emph{Shadow-model} attacks established the threat: by training
reference models on data drawn from the same distribution, an adversary learns to
distinguish members from non-members from the target model's
outputs~\cite{shokri2017membership}. Yeom~et~al.\ tied attack success to overfitting and
gave the simplest practical baseline (thresholding the per-example loss) together with the
\emph{membership advantage} (TPR${-}$FPR) figure of merit~\cite{yeom2018privacy}.
Carlini~et~al.'s \emph{Likelihood Ratio Attack} (LiRA) then reframed MIA from first
principles as a per-example hypothesis test calibrated with shadow models, and, central
to our methodology, argued that average-case AUC is the wrong yardstick for a privacy
threat: an attack matters if it identifies \emph{some} members with very few false
accusations, so the right report is TPR at a low, fixed FPR on a log-scale ROC
curve~\cite{carlini2022lira}. Shadow-model calibration, however, is infeasible at
Pile/Pythia scale (it requires training many models on the training distribution), so we
adopt LiRA's \emph{metric} but not its \emph{attack}.

For pre-trained LLMs, the field moved to \emph{reference-free} likelihood signals that
need no shadow models. Min-K\% Prob averages the log-probabilities of a sequence's
lowest-probability $k\%$ of tokens, on the hypothesis that members lack
high-surprise outlier tokens~\cite{shi2024detecting}; Min-K\%++ sharpens this by
$z$-scoring each token against the \emph{full} next-token distribution before averaging,
detecting that the target token sits at a local maximum of the modeled
distribution~\cite{zhang2025minkpp}. A parallel reference-free line, neighbourhood
comparison, calibrates a sample's score against synthetically generated neighbour texts
instead of a reference model~\cite{mattern2023neighbourhood}; we treat it as a related
approach we do not evaluate, since it needs many extra masked-LM forward passes per
example and, in the regime below, underperforms. The reality check on this whole line is
the MIMIR study: a large-scale audit on Pythia ($160$M--$12$B) and The Pile with
controlled member/non-member splits finds that these attacks barely exceed chance
(AUC~$\approx 0.5$--$0.6$), that LLMs see their corpus for too few epochs over too large
a dataset to memorize in the way classical MIA assumes, and that apparent successes
frequently reflect a temporal or topical \emph{distribution shift} between the splits
rather than membership~\cite{duan2024mia}. This finding defines our honesty constraint:
we do not claim to beat these numbers; we ask whether the weak signal that remains still
predicts leakage.

\subsection{Differential privacy as the defense direction}
\label{sec:dp}
The standard principled mitigation for training-data leakage is differential privacy.
DP-SGD bounds any single example's influence on the trained model by clipping per-example
gradients and adding calibrated noise, with privacy accounted via the moments
accountant~\cite{abadi2016deep}. Applied to language models, DP fine-tuning can retain
much of the utility of non-private training, particularly with large pre-trained
backbones~\cite{li2022dpllm} and parameter-efficient adaptation~\cite{yu2022dpfinetuning}.
DP bounds memorization and thereby the leakage we measure, but at a privacy--utility cost
and, crucially for us, it must be applied \emph{at training time}; it is a defense for
model producers, not a detector available to an auditor of an already-released model. We
therefore position DP as the mitigation our threat model motivates, and do not implement
it (we train no models).

\subsection{Existing detection techniques}
We describe the techniques we implement and compare; the comparative evaluation and
the access requirements appear in Section~\ref{sec:eval}. All operate without any novel
detector of our own, our contribution is their security-framed, ground-truth
evaluation, not a new method.
\begin{itemize}
 \item \textbf{$n$-gram / substring overlap.} Flag a benchmark item that shares an
 $N$-gram with the corpus~\cite{brown2020gpt3}. Requires corpus access; misses
 paraphrased and semantic contamination.
 \item \textbf{Loss / perplexity thresholding.} The mandatory membership-inference
 baseline: members exhibit lower loss, with attack success tied to
 overfitting~\cite{yeom2018privacy}.
 \item \textbf{Min-K\% Prob.} Average the log-probabilities of the lowest-probability
 $k\%$ of tokens; reference-free and logprob-only~\cite{shi2024detecting}.
 \item \textbf{Min-K\%++.} Normalizes each token's log-probability against the full
 next-token distribution before the bottom-$k\%$ average, the current state of the art
 among reference-free detectors~\cite{zhang2025minkpp}.
 \item \textbf{zlib-entropy ratio.} Calibrate model perplexity by the zlib-compressed
 size of the text, controlling for intrinsic compressibility/frequency~\cite{carlini2021extracting}.
 \item \textbf{Permutation / exchangeability test.} At the \emph{benchmark} level rather
 than per item, score each ordering of a benchmark's examples by the log-likelihood of
 their concatenation and compare the canonical (published) order against random
 shufflings; a model trained on the benchmark in canonical order favours it beyond
 chance, yielding a provable, FPR-controlled contamination
 certificate~\cite{oren2024proving}.
\end{itemize}
We additionally note two techniques we describe but \emph{do not} evaluate, since our
ground-truth, logit-access setting makes likelihood-based detectors stronger and cleaner:
\emph{guided prompting}, which prompts a model with dataset metadata and a partial
instance and tests for verbatim completion~\cite{golchin2024timetravel}, a black-box
signal aimed at closed models; and the reference-free \emph{neighbourhood} and
shadow-model \emph{reference} attacks discussed in
Section~\ref{sec:mia-lineage}~\cite{mattern2023neighbourhood,shokri2017membership}.

\subsection{Limitations of existing detection, and our positioning}
Two limitations frame our contribution. First, \emph{detection is fragile to the
transformation}: string-matching misses paraphrased and semantic
contamination~\cite{ippolito2023verbatim}, and likelihood-based membership inference is
known to barely exceed chance on pre-trained LLMs evaluated under controlled ground
truth, because the corpora are seen for few epochs and member/non-member boundaries are
fuzzy~\cite{duan2024mia}. Second, \emph{evaluation conventions matter}: average-case
AUC or accuracy can mask whether an attack confidently identifies any members, so the
security-appropriate report is true-positive rate at low false-positive rate with
log-scale ROC~\cite{carlini2022lira}. We therefore do not claim a stronger detector.
We ask a different, security-relevant question: \emph{even where contamination signal
is weak, does it predict concrete privacy leakage?} We answer it with ground-truth
membership on the Pythia suite~\cite{biderman2023pythia} trained on the public
Pile~\cite{gao2020pile}, under the low-FPR protocol, with explicit controls for the
frequency, duplication, and temporal confounds that prior work identifies.

\subsection{Closest prior work, and how we differ}
\label{sec:closest}
Three recent works reach conclusions adjacent to ours, and we are careful to position
against them rather than overclaim. Al Sahili et al.~\cite{alsahili2025effectiveness}
reach a compatible conclusion for targeted extraction, that ``complex MIA techniques
yield only marginal improvements over simple likelihood-based ranking'', but they
establish it through aggregate \emph{ranking-precision} comparisons and an AdaBoost
ensemble over MIA features, reporting \emph{marginal gains} rather than testing for
independent signal. In contrast, we run a pre-registered \emph{partial correlation
controlling for raw per-item loss}, which lets us state the stronger, calibrated claim
that the reference-free detectors contribute \emph{zero or negative} residual predictive
value once loss is partialled out. Hayes et al.~\cite{hayes2025strong} likewise
``observe no correlation with MIA success'' for extraction and conclude the ``two privacy
attacks may capture different signals,'' but their evidence is a \emph{direct, zero-order}
correlation between a reference-model attack (LiRA) and extraction. We differ on both
method and object: we \emph{partial out per-item loss} rather than correlating directly,
and we target the reference-free \emph{calibrated} detectors (Min-K\%, Min-K\%++, zlib)
that the contamination-detection literature actually deploys, showing the divergence
persists as a controlled mediation result. Independently, Chen et
al.~\cite{chen2025statistical} find for the \emph{membership} task that the few detectors
numerically above the loss baseline (Min-K\%, Min-K\%++, ReCaLL) do not beat it robustly
once random-seed variance is accounted for, and that performance is domain-dependent
(code-like, low-token-diversity domains such as GitHub and StackExchange behave
differently from Wikipedia and FreeLaw); we revisit this domain dependence for the
\emph{extraction} outcome in our per-domain analysis (Section~\ref{sec:eval}), noting it
is a distinct axis from their membership-AUC result. Finally, blind-baseline and SoK
critiques~\cite{das2024blind,meeus2025sok} show that post-hoc member/non-member splits can
make detector ``success'' an artifact of distribution shift; our use of ground-truth Pile
membership (no post-hoc split) is precisely the design discipline they call for.

\begin{table}[t]
\centering
\footnotesize
\setlength{\tabcolsep}{4pt}
\begin{tabular}{@{}p{2.1cm}p{1.6cm}p{1.9cm}p{2.0cm}p{2.4cm}@{}}
\toprule
\textbf{Study} & \textbf{Outcome} & \textbf{Detectors} & \textbf{Statistical method} & \textbf{Conclusion} \\
\midrule
Shi'24; Zhang'25~\cite{shi2024detecting,zhang2025minkpp} & membership & reference-free (Min-K\%/++) & AUC / TPR@FPR & detector raises membership AUC \\
Duan'24 (MIMIR)~\cite{duan2024mia} & membership & ref-free + reference & AUC on ground truth & MIAs $\approx$ chance on LLMs \\
Carlini'22 (LiRA)~\cite{carlini2022lira} & membership & shadow/reference & TPR at low FPR & strong only with shadow models \\
Chen'25~\cite{chen2025statistical} & membership & reference-free & seed-variance testing vs loss & not robustly beyond loss \\
Hayes'25~\cite{hayes2025strong} & membership \& extraction & LiRA (reference) & direct (zero-order) correlation & MIA $\neq$ extraction \\
Al Sahili'25~\cite{alsahili2025effectiveness} & extraction (targeted) & ref-free + AdaBoost & ranking precision; ensemble & marginal gains over likelihood \\
\textbf{This work} & \textbf{extraction} & \textbf{ref-free calibrated} & \textbf{partial corr.\ + mediation (control loss)} & \textbf{zero/negative residual beyond loss} \\
\bottomrule
\end{tabular}
\caption{Where this work sits. To our knowledge it is the only study that pairs a per-item
\emph{extraction} outcome with a \emph{partial-correlation/mediation} control for raw loss
on \emph{calibrated reference-free} detectors, yielding a quantified zero/negative marginal.}
\label{tab:closest}
\end{table}
